\documentclass[10pt,letterpaper]{article}
\usepackage[top=0.85in,left=2.75in,footskip=0.75in]{geometry}

% amsmath and amssymb packages, useful for mathematical formulas and symbols
\usepackage{amsmath,amssymb}

% Use adjustwidth environment to exceed column width (see example table in text)
\usepackage{changepage}

% textcomp package and marvosym package for additional characters
\usepackage{textcomp,marvosym}

% cite package, to clean up citations in the main text. Do not remove.
\usepackage{cite}

% Use nameref to cite supporting information files (see Supporting Information section for more info)
\usepackage{nameref,hyperref}

% line numbers
% \usepackage[right]{lineno}

% ligatures disabled
\usepackage[nopatch=eqnum]{microtype}
\DisableLigatures[f]{encoding = *, family = * }

% color can be used to apply background shading to table cells only
\usepackage[table]{xcolor}

% array package and thick rules for tables
\usepackage{array}

% create "+" rule type for thick vertical lines
\newcolumntype{+}{!{\vrule width 2pt}}

% create \thickcline for thick horizontal lines of variable length
\newlength\savedwidth
\newcommand\thickcline[1]{%
  \noalign{\global\savedwidth\arrayrulewidth\global\arrayrulewidth 2pt}%
  \cline{#1}%
  \noalign{\vskip\arrayrulewidth}%
  \noalign{\global\arrayrulewidth\savedwidth}%
}

% \thickhline command for thick horizontal lines that span the table
\newcommand\thickhline{\noalign{\global\savedwidth\arrayrulewidth\global\arrayrulewidth 2pt}%
\hline
\noalign{\global\arrayrulewidth\savedwidth}}

% Remove comment for double spacing
%\usepackage{setspace} 
%\doublespacing

% Text layout
\raggedright
\setlength{\parindent}{0.5cm}
\textwidth 5.25in 
\textheight 8.75in

% Bold the 'Figure #' in the caption and separate it from the title/caption with a period
% Captions will be left justified
\usepackage[aboveskip=1pt,labelfont=bf,labelsep=period,justification=raggedright,singlelinecheck=off]{caption}
\renewcommand{\figurename}{Fig}

% Use the PLoS provided BiBTeX style
\bibliographystyle{plos2015}

% Remove brackets from numbering in List of References
\makeatletter
\renewcommand{\@biblabel}[1]{\quad#1.}
\makeatother


% Header and Footer with logo
\usepackage{lastpage,fancyhdr,graphicx}
\usepackage{epstopdf}
%\pagestyle{myheadings}
\pagestyle{fancy}
\fancyhf{}
%\setlength{\headheight}{27.023pt}
%\lhead{\includegraphics[width=2.0in]{PLOS-submission.eps}}
\rfoot{\thepage/\pageref{LastPage}}
\renewcommand{\headrulewidth}{0pt}
\renewcommand{\footrule}{\hrule height 2pt \vspace{2mm}}
\fancyheadoffset[L]{2.25in}
\fancyfootoffset[L]{2.25in}
\lfoot{\today}

\newcommand{\orcid}[1]{\href{https://orcid.org/#1}{#1}}


\begin{document}
\vspace*{0.2in}

\begin{flushleft}
{\Large
\textbf\newline{10 quick tips for making your software outlive your job}
}
\newline
\\
% Clare Dillon\textsuperscript{2},
Priyanka Ojha\textsuperscript{3},
Georg Link\textsuperscript{5},
Mala Kumar\textsuperscript{6},
Eman Abdullah AlOmar\textsuperscript{9},
Mohamed Wiem Mkaouer\textsuperscript{10},
David Lippert\textsuperscript{11},
Daniel R.\ McCloy\textsuperscript{12},
Bill Branan\textsuperscript{13},
Chang Liao\textsuperscript{16},
Christoph Treude\textsuperscript{17},
Tobias Augspurger\textsuperscript{18},
Sylwester Arabas\textsuperscript{19},
Ethan P.\ White\textsuperscript{20},
Fang Liu\textsuperscript{21},
Geoffrey Lentner\textsuperscript{22},
David Eyers\textsuperscript{23},
Elena Findley-de Regt\textsuperscript{25},
Paola Corrales\textsuperscript{27},
Pieter Huybrechts\textsuperscript{28},
Phani Velicheti\textsuperscript{29},
Daniel Morillo-Cuadrado\textsuperscript{30},
Tommy Guy\textsuperscript{31},
Jan Ainali\textsuperscript{32},
Jonathan Guyer\textsuperscript{33}
\\
\bigskip
\textbf{2} CURIOSS, Lero, University of Galway, Ireland, \orcid{0009-0008-6205-0296}\\
\textbf{3} \orcid{0000-0002-6844-6493}\\
\textbf{5} Bitergia, \orcid{0000-0001-6769-7867}\\
\textbf{6} unknown\\
\textbf{9} \orcid{0000-0003-1800-9268}\\
\textbf{10} \orcid{0000-0001-6010-7561}\\
\textbf{11} \orcid{0009-0003-6444-9595}\\
\textbf{13} \orcid{0000-0002-4735-6624}\\
\textbf{16} \orcid{0000-0002-7348-8858}\\
\textbf{17} \orcid{0000-0002-6919-2149}\\
\textbf{18} Unknown\\
\textbf{19} AGH University of Krakow, Poland \orcid{0000-0003-2361-0082}\\
\textbf{20} Department of Wildlife Ecology and Conservation, University of Florida, Gainesville, FL, USA \orcid{0000-0001-6728-7745}\\
\textbf{21} \orcid{0000-0002-3383-2191}\\
\textbf{22} \orcid{0000-0001-9314-0683}\\
\textbf{23} \orcid{0000-0002-7284-8006}\\
\textbf{25} Unknown\\
\textbf{27} \orcid{0000-0003-1923-9129}\\
\textbf{28} Research Institute for Nature and Forest (INBO), Brussels, Belgium, \orcid{0000-0002-6658-6062}\\
\textbf{29} \orcid{0009-0004-2580-3624}\\
\textbf{30} Universidad Nacional de Educación a Distancia (UNED), Spain, \orcid{0000-0003-3021-3878}\\
\textbf{31} \orcid{0009-0003-3652-5036}\\
\textbf{32} Open By Default, \orcid{0000-0001-8747-1670}\\
\textbf{33} \orcid{0000-0002-1407-6589}\\

Organized by Richard Littauer: see \texttt{contributors.tex} for all contributors to date.
\end{flushleft}

\section*{Abstract}

We present ten tips to help researchers ensure that
the software they have built will continue to be usable
after they have left their present job.

\linenumbers

\section*{Introduction}

Academic life is often precarious,
particularly for those who work on research software,
an asset that is essential to modern research \cite{Pearson2025}
but often underfunded, undercited, and overlooked \cite{Carver2022}.
The situation has recently worsened dramatically due to changes in the political climate:
researchers in all disciplines face the loss of funding or jobs,
a smaller pool of incoming students,
and institutional policies that prioritize commercial returns.
Researchers in non-academic organizations
such as government labs, non-governmental organizations (NGOs), and non-profits are also impacted \cite{Woodward2025}
as national research strategies are upended
and associated public funding redirected or cut entirely \cite{Nature2025}.

In this environment it is essential that research software projects be \emph{resilient}.
Code should outlast situations the author changing workplaces or careers,
partnerships failing,
a lab closing down,
or tooling and digital infrastructure becoming unavailable.
We therefore present ten tips to ensure your software remains accessible and usable by others
even if you are no longer able to work on it.
You do not need to adopt all of them to make things better;
instead,
as in all emergencies,
you should do what you can,
where you are,
with what you have.

This guide focuses on software.
Other guides already exist for data \cite{Perkel2023},
and archives such as \href{https://zenodo.org/}{Zenodo},
\href{https://figshare.com/}{FigShare},
the \href{https://eotarchive.org/}{End Of Term Archive}
or the \href{http://archive.org/}{Internet Archive}
are now widely used to ensure long-term availability.
However,
we have not found any guide to making research software resilient
or to \emph{sunsetting} it (i.e., phasing it out gracefully).
Even if your job is not in danger,
these tips will help others access and use your code,
contribute to your project,
and to reproduce and cite your work.
As such,
our recommendations
are part of the global push for more equitable assessment of research contributions
by initiatives like \href{https://sfdora.org/}{DORA},
\href{https://coara.eu/}{coARA},
and \href{https://adore.software/}{ADORE.software}.

\section*{Tip 1: Consider your threat model}

When making plans, it's useful to know what you're planning \emph{for}.
Being explicit about \textbf{threat models} helps you prioritize,
build consensus with colleagues,
and check whether you have forgotten something important.
As with all disasters,
making those plans before you need them helps reduce the odds of them being needed,
as planning may help you identify risks you can eliminate.

\begin{enumerate}
\item
  \textbf{Individual threats} affect one or a few members of your team,
  such as a foreign student having their visa revoked without notice
  or a contributor taking extended leave.
  Especially for software under the maintenance of a single author,
  individual threats can also come from more mundane career or life changes.
  The most common way to prepare for this is to require everyone to document their work,
  but that rarely works in practice:
  \begin{enumerate}
  \item
    The hours spent writing those descriptions are hours \emph{not} spent doing research,
    so people will always short-change the former to focus on the latter.
  \item
    People invariably fail to write down the ``obvious'' parts of their work
    that are anything but obvious to the next person.
  \end{enumerate}
  Tip~8 explores alternatives,
  particularly ones that can be put in place on short notice.

\item
  \textbf{Leadership threats} are individual threats that affect the project's leader,
  such as the leader being doxxed or targeted personally in the media because of their work.
  Tip~3 discusses ways to plan for this.

\item
  \textbf{Institutional threats} affect large groups at once,
  such as your department being shut down
  or your entire field having funding cut.
  These events affect so many people at the same time
  that the rest of the community can't absorb them.
  Regional and national governments plans for disasters like these
  by having evacuation plans to get victims to safe(r) places
  and corollary plans for putting beds in high school gyms
  and flying in food and emergency medical personnel to help people when they arrive.
  At the time of writing,
  universities and researchers' professional societies have not started to do the equivalent.

\item
  \textbf{Global threats} affect everyone,
  not just researchers.
  For example,
  there is no technical or legal obstacle to the US government
  requiring American companies to charge for video conferencing calls
  involving participants outside the United States.
  Similar levies on email,
  file storage,
  and other online services would paradoxically have \emph{less} long-term impact on research
  than the targeted threats described above,
  as they would force national governments to find effective remedies quickly.

\end{enumerate}

\section*{Tip 2: Stay within the law}

Ensuring your work remains usable is not worth putting yourself at legal risk.
Before following any of the tips below,
make sure you have a legal right to do so.
Most institutions and journals now have policies for licensing code and/or releasing it publicly \cite{Katz2018,Ham2019}.
Funding agencies at various levels also often have policies
(e.g.,
EU\footnote{\url{https://commission.europa.eu/about/departments-and-executive-agencies/digital-services/open-source-software-strategy_en}}
and NASA\footnote{\url{https://science.nasa.gov/open-science/nasa-open-science-funding-opportunities}})
which may or may not align with those of institutions and journals,
and institutional policies may place restrictions on what you are allowed to say publicly about your situation.
Our first actionable tip is therefore to find out what those policies are,
e.g.,
whether you need explicitly permission to publish your software,
and if so,
from whom.

If you are unsure whether there is a policy or not,
ask your direct manager,
the person who pays you,
or offices like the research office,
the library,
or the tech transfer office.
If your institution has a dedicated
\href{https://sustainoss.org/academic-map/universities/index.html}{Open Source Program Office} (OSPO),
ask them
or connect to networks like \href{https://curioss.org}{CURIOSS},
the \href{https://ospo-alliance.org/}{OSPO Alliance}, 
or the \href{https://todogroup.org/}{TODO Group}
that advocate for OSPOs.
If you work with multiple institutions,
look at your contract and the bilateral agreements around products or deliverables.

Once you know whether your institution has a policy on licensing code,
get whatever sign-offs you need immediately.
Reach out to colleagues who can review your code if that is necessary for publishing it,
and review theirs if they ask.
If it appears that there is no formal sign-off process,
send an email to someone in authority
(e.g., the chair of your department or your grant officer)
saying explicitly that you believe this to be the case,
and copy that email to an account you will be able to access
after leaving your institution.

Do not assume that if you had permission before,
you have it now or will have it in the future.
Policies are changing rapidly,
and you may find yourself locked into one that no longer allows you to do what you want.
Also consider that the person you report to may be replaced with one who knows you less well or is less sympathetic,
so acting now may be easier than acting later.

Finally,
if you have your next job lined up,
ask about their policies
and make sure that your right to share your work is written into your contract.
Making your code comply with policy after the fact is riskier and more time-consuming than doing so early,
and sometimes not possible.

\section*{Tip 3: Decide if you are ending, pausing, and passing on}

There are thousands of books in print about how to start a business,
but only a handful describe how to hand one over or wind one down,
and most of those focus on succession within family-owned businesses.
Guides to running a laboratory such as \cite{Barker2010,Cohen2018} say little more,
so such wisdom as we have is passed on person-to-person if it is passed on at all.

The first step is to decide if work on the project is ending,
being suspended temporarily (for some version of ``temporarily''),
or if you are handing it on to someone else.
If the project is ending or going on hiatus,
focus on making what you have accomplished findable and citable
as discussed in the tips that follow.
To help you prioritize,
write up the goals you initially had for this project
and ask how well they were served by what you actually did.
Doing this may actually help you realize that the project doesn't have to end,
but can instead be folded into someone else's work,
even if you are no longer involved.

If instead you hope to hand the project on to someone else,
you can either choose a successor
(who will typically be someone already working on the project)
or ask your peers for volunteers or recommendations.
Writing a \emph{succession plan} that describes what the work will actually entail
will help in both cases.
What will your successor have to learn and do during the transition?
What will they be responsible for once they are in charge?
And crucially,
how much work is running the project actually going to be?
Please be honest rather than optimistic:
finding out after you have agreed to take something on
that it is much larger than you realized
can strain both professional and personal relationships.

If the project is large,
and you have enough time,
consider working with a fiscal sponsor such as
\href{https://numfocus.org/}{NumFOCUS},
\href{https://sfconservancy.org/}{Software Freedom Conservancy},
\href{https://www.eclipse.org/}{the Eclipse Foundation},
\href{https://oscollective.org/}{Open Source Collective},
or other \href{https://sustainoss.org/academic-map/organizations/index.html}{foundations and fiscal hosts}.
Through fiscal hosting,
your project can take donations from its community to fund continued maintenance and other costs.
Universities and governments are generally poorly suited for this due to high overheads and procurement costs;
fiscal sponsors may take an order of magnitude less.
In addition,
a fiscal sponsor puts the software in a neutral home
where it is no longer owned by a single research organization,
which helps insulate it from shifts in policy.

\section*{Tip 4: Choose an open license}

Making code publicly available does not ensure that other people can use it:
you maintain all copyright unless you explicitly include a license,
so they may still need explicit permission
or may simply be unsure whether they do or not.
The Open Source Initiative (OSI)
maintains a list of \href{https://opensource.org/licenses}{open source licenses} it has approved;
their implications are widely understood,
so choosing one of those will make your project more understandable to others.
Please do \emph{not} try to write your own license,
as it will have the opposite effect.

Broadly speaking,
the MIT, Apache 2.0, or BSD-3-Clause licenses place the fewest restrictions on re-use.
The GPLv3 and AGPL-3.0 licenses are \emph{copyleft} licenses
that require people to share their modifications to the software with the community \cite{Morin2012}.
These licenses can prevent companies from taking advantage of your work without giving anything in return,
but institutions discourage or disallow the use of copyleft licenses
out of fear that they will constrain commercialization.

More broadly,
\href{https://creativecommons.org/}{Creative Commons} licenses can be used for documentation or research reports,
while \href{https://ethicalsource.dev/}{ethical licenses}
or licenses based on the \href{https://blueoakcouncil.org/license/1.0.0}{Blue Oak Model License} may also be appropriate.
For guidance,
go to \href{http://choosealicense.com}{choosealicense.com}
or refer to \cite{Fogel2020,Fortunato2021};
the discussion in \cite{Gomes2022} about why people \emph{don't} share is also helpful.
Whatever license you choose,
place its text in a file called \texttt{LICENSE}, \texttt{LICENSE.txt}, or \texttt{LICENSE.md}
in the root directory of your project,
as this is where people and automated tools alike will look for it.

\section*{Tip 5: Save everything in multiple places}

Now that you're legally able to share your code, remember LOCKSS: Lots of Copies Keep Stuff Safe.
But copies are not enough:
the threats to your project are political as well as technological,
so you should ensure that loss of institutional support
(or worse, that institution turning on you)
does not mean loss of access.

\begin{quote}
  \noindent
  \textbf{The Road Not Taken}

  \cite{Langille2010} described a file-sharing system for scientific data
  based on the BitTorrent protocol.
  It failed to find wide adoption because of the (quite reasonable) association in many institutions' collective minds
  between BitTorrent and illegal downloading.
  At the time of writing,
  though,
  many individuals and groups are turning to BitTorrent to share datasets that are at risk of disappearing
  precisely because of the protocol's resilience.
\end{quote}

\href{https://github.com/}{GitHub},
\href{https://gitlab.com}{GitLab},
and \href{https://bitbucket.org/}{BitBucket}
are social coding platforms centered around Git,
a widely-used open source version control tool.
However,
they are all commercial entities based in the same legal jurisdiction,
which means they are vulnerable to \emph{correlated threats}:
trouble with any of them may mean trouble with all of them.
(\cite{Tamburri2020} is a sobering account of what can happen when key infrastructure comes under threat.)
\href{https://codeberg.org/}{Codeberg},
while currently much smaller,
is a non-profit based in a different jurisdiction;
consider mirroring your work there or using it as your primary host.

When you create projects on hosted services,
use a team account as the project owner rather than a personal account.
Doing this makes it easier to give other people permission to manage the project,
and as noted in Tip~6,
a project with multiple owners from different institutions
is harder for any one institution to lock down.
For the same reason,
do not rely solely on logins or email addresses associated with your institution:
instead,
ask yourself if you will still have access to the project if you lose that identity,
and add an identity you personally control to the team that owns the project.
You can also add your \href{https://orcid.org/}{ORCID} to the project's metadata
so that the project links to a profile that you control
even when your contact address changes.

\begin{quote}
  \noindent
  \textbf{What's In a Project?}

  When deciding what to store where,
  remember that your version control repository only stores part of what makes up your project.
  For example,
  GitHub issues and wikis are displayed in the browser but live in GitHub's database;
  you can use the GitHub API to download their content,
  but (a) the result is intended for consumption by machines, not people,
  and (b) if you are doing this on short notice,
  the odds are high that other people are trying to do it at the same time.
  GitHub's servers may not be able to manage that load when you need them to,
  so use tools like \href{https://github.com/jlord/offline-issues}{offline-issues}
  to save hosted content as plain-text files on a regular basis.
\end{quote}

Version control isn't the only way to create and save copies of your work.
\href{https://www.softwareheritage.org/how-to-archive-reference-code/}{Software Heritage} archives software from multiple forges;
you can also snapshot the current state of repositories in a compressed archive file
(e.g., \texttt{.zip} or \texttt{.tar.gz})
and deposit those copies with the \href{https://osf.io/}{Open Science Foundation} (OSF),
\href{https://zenodo.org/}{Zenodo},
or \href{https://figshare.com/}{figshare}.
(Zenodo even integrates with GitHub so that tagging a release on GitHub
automatically triggers deposit of a new archive on Zenodo,
but again,
this automation is only useful as long as both ends are accessible.)
Institutional,
\href{https://amt.coretrustseal.org/certificates/}{national},
or \href{https://safeguar.de/}{international} data repositories also enhance longevity.

Storing copies on someone else's computer isn't the only option:
you can (and should) make copies of your projects on computers that you own.
Again,
this is a place where colleagues can help:
ask them to make copies on their computers in exchange for you making copies of theirs on yours.
Whatever you do,
document it \href{https://ropensci.org/blog/2022/03/22/safeguards-and-backups-for-github-organizations/}{like rOpenSci has}.

Similarly,
if you distribute your code as a package on \href{https://pypi.org/}{PyPI},
\href{https://anaconda.org/anaconda/conda}{Conda},
or \href{https://cran.r-project.org/}{CRAN},
that package can contain all of the source code.
Doing this also fosters community collaboration

\section*{Tip 6: Encourage community adoption}

Publishing your code is not the same as publicizing it.
the more people who know about and rely on your project,
the greater the odds that someone will help keep it alive,
whether by working on directly or supporting you to do so.
A simple way to help people get involved is to label entry-level issues
(e.g., to use the ``good first issue'' label on GitHub)
so that people can find a place to start \cite{Steinmacher2015}.

If you haven't named your project yet,
choose a name that doesn't hint at direct and exclusive affiliation to a single institution.
This can help if the code needs to be relocated,
and also makes contributors from other institutions feel welcome.

Announcing your project on mailing lists, forums, or social media,
and talking about it at conferences are necessary but not sufficient:
everyone else is doing this too,
so your project will almost certainly get lost in the noise.
\cite{Kuchner2011} is a good guide to what you can do beyond the obvious.
A short video showing how to use the software
or a slide deck that others can incorporate into their lectures
will increase uptake by lowering the cost of adoption.
Similarly,
a one-page website that opens with an elevator pitch
explaining who the software is for
and how it will make their lives better
is much more likely to lead to that crucial ``second glance''
than a list of publications.

Disseminating your code via tutorials at conferences is another effective strategy.
Tutorials are an opportunity for you,
as an author,
to describe your work in depth
and (more importantly) convince your audience to use it.
Networking through such events will enable you to build relationships with
the people most likely to care about your project,
and is the best way to ensure that you land well after you leave your current position.

\begin{quote}
  \noindent
  \textbf{One to Many}

  If you can, have at least one person outside your organization commit to your code.
  Such community contributions lead to joint ownership of intellectual property (IP),
  which makes it harder for any single institution to lock down your work.
  You can reinforce this by adding their name to the copyright statement in your license.
\end{quote}

Finally,
remember that it isn't all about \emph{your} project.
If you ask others to help maintain your code,
be prepared to give them something in return:
a letter of reference,
help testing or documenting their project,
or co-authorship credit on the project and associated publications.

\section*{Tip 7: Do what you're supposed to (if you have time)}

Research software is increasingly recognized as a citable object \cite{Smith2016,Katz2021,Garijo2024},
and numerous books and papers describe how to organize and run a research software project.
Rather than repeating what is in
\cite{Sandve2013,Wilson2014,Lee2018a,Dryden2019,Akhmerov2020,Chue2021,Lees2022,Druskat2023,Akhmerov2023,Kumar2023,Struck2023,Reina2024},
our sixth tip is a prioritized checklist of things to do
\emph{if you have time}.
If you do not,
please consult the next tip.

\begin{enumerate}

\item
  Create a DOI for each release of your software
  to ensure citability and retrievability even if the project repository becomes unavailable,
  and add a \href{https://citation-file-format.github.io/}{Citation.CFF}
  or \href{https://codemeta.github.io}{codemeta.json} file to your code \cite{Druskat2021}
  containing citation metadata.
  Tools such as \href{https://citation-file-format.github.io/cff-initializer-javascript/}{cffinit}
  \href{https://codemeta.github.io/codemeta-generator/}{codemeta generator},
  and \href{https://github.com/citation-file-format/citation-file-format/blob/main/README.md\#tools-to-work-with-citationcff-files-wrench}{others}
  can help you do this.

\item
  Have your code peer-reviewed and published in a venue such as
  the \href{https://joss.theoj.org/}{Journal of Open Source Software},
  the \href{https://openresearchsoftware.metajnl.com/}{Journal of Open Research Software},
  the \href{http://www.jstatsoft.org}{Journal of Statistical Software},
  the \href{https://journal.r-project.org/}{R Journal},
  or \href{https://ropensci.org/}{rOpenSci}.

\item
  Document how to \emph{use} your software so that others can pick it up and use it if you are not available.
  In particular,
  add a \texttt{README.md} at the root of your repository
  that explains the purpose, setup, and use of your work,
  and add a tutorial showing how to use the major components of the code
  \cite{Lee2018b,Huybrechts2024,Littauer2025,Katz2025,Turing2025}.

\item
  Document any datasets your project depends on,
  and use CSV, JSON, Parquet, HDF5, or other widely-recognized formats rather than proprietary formats.
  Include the data in your project if possible;
  if that is not possible for legal reasons or because of the data's size,
  include URLs with data version IDs (or the date of the last download) in your documentation.

\item
  Describe your project's dependencies in a machine-readable way
  using Python's \href{https://packaging.python.org/en/latest/guides/writing-pyproject-toml/}{pyproject.toml} file,
  an npm \href{https://docs.npmjs.com/cli/v10/configuring-npm/package-json?v=true}{package.json} file,
  the \texttt{DESCRIPTION} and \texttt{NAMESPACE} files of an R package,
  or whatever else is standard for your language.

\item
  Write at least a few tests for your software
  so that people who want to use it (or contribute to it)
  can tell if it's working in their environment.
  These do not have to be formal unit tests \cite{Irving2021}:
  as \cite{Taschuk2017} explains,
  even a handful of ``this input should produce this output'' checks
  can save a lot of frustration.

\item
  Document how to \emph{contribute} to your software,
  e.g.,
  how to set up a development environment,
  run tests,
  and add new features.
  Such a guide leads to an increase in contributions \cite{Sholler2019},
  as does an explanation of the project's governance,
  i.e.,
  how decisions are made and who gets a voice in making them.

\end{enumerate}

\section*{Tip 8: Do what you can (if you have to act quickly)}

A paramedic is not a doctor who cuts corners,
but rather someone who has been trained to work subject to very different constraints.
\cite{Wilson2017} provides some guidance for ``good enough'' practices in research computing;
other tips in this paper are in this vein,
such as creating \texttt{.zip} archives of your project as an alternative to using version control.
A few others include:

\begin{enumerate}

\item
  Explanations of what the software does and how to use it
  take priority over documentation of its internals
  because (a) someone who knows the former will be better able to figure out the latter
  and (b) AI code assistants can do a reasonable job of helping newcomers navigate an unfamiliar code base
  if they are given enough context.

\item
  Similarly,
  a short screencast showing how you use your software
  is often faster to create than pages of documentation,
  and may actually be more useful provided provided closed captions are attached
  to make the content searchable.
  Again,
  AI tools can help create those captions,
  though some editing is invariably required to clean up mistranslation of technical jargon.

\item
  If you do not have time to bring your project in line with best practices,
  try instead to clean it up by removing files that are no longer used,
  deleting sections of your setup instructions that no longer apply,
  closing bug reports that are no longer relevant,
  and so on.
  Eliminating clutter in this way will make what's useful easier to find.
  
   \item
  Briefly document what happened.
  If you're the main developer working on a grant, and you've moved on to something else, say that in the README.
  If the funding changed, say that the work is unfunded.
  This information is often the most useful for users, but the hardest to find. Saying it up front can help others more readily understand how your work is positioned.

\end{enumerate}

\section*{Tip 9: Talk about what you're doing}

All of the work we have described is extra labor.
You won't want to do it,
particularly not when so many other things demand attention as well.
Make that work meaningful by helping others.
Tell your colleagues what you have learned about your institution's legal requirements,
about navigating its bureaucracy,
and about the least hurtful way to break the news to your students, collaborators, and users.
Share your own tips on social media
(preferably Mastodon rather than X or Bluesky,
as the latter two are single points of institutional failure),
and share things that \emph{haven't} worked as well
so that other people can avoid stepping on the same landmines.

Most importantly,
take the time to grieve.
Having years of work come to an abrupt end will take a toll on your mental health;
many of your colleagues will be suffering as well,
and those who haven't been hit (yet) may be affected by survivors' guilt.
Be angry,
be sad,
but remember that in the aftermath of a disaster---natural or man-made---people
often become altruistic, resourceful, and brave.

\section*{Tip 10: Organize.}

There is little point surviving today's flood if you drown in tomorrow's.
The best strategy is to take preventive action,
and the best way to do that in a research context
is to become active in professional associations
and push them to take meaningful action.

Many readers of this paper will find the thought of doing this uncomfortable.
As \cite{Putnam2020} points out,
voluntary organizations have seen a slow but steady decline in membership over several decades.
While there are many contributing factors,
the net result is that professionals aged 30--60
are less likely to be involved in civil society and the political process than the generation before them.
This disengagement effectively hands power to extremists and special interest groups
and enables them to seize control of civil institutions in order to advance their agendas \cite{BuenoDeMesquita2022}.
To see the impact,
ask yourself who is more likely to run for a seat on your local school board:
a professor of epidemiology or an anti-vax conspiracy theorist?

We have written this paper because we believe scientific inquiry is a public good and worth fighting for.
The present climate requires us to focus on preserving what we can,
but playing defense is just a way to lose more slowly.

\bibliography{code-rules}

\end{document}
